\subsection{Application Domains: Awareness in Practice}
\label{subsec:11-16-applications}

%%%%%%%%%%%%%%%%%%%%%%%%%%%%%%%%%%%%%%%%%%%%%%%%%%%%%%%%%%%%%%%%%%%%%%%%%%%%%%%
% SECTION 11.16: Application Domains
% All 9 cases worked through: Intuitive + Formal + NUMERICAL CALCULATION
% Updated: January 2026 with explicit U_eff computations
%%%%%%%%%%%%%%%%%%%%%%%%%%%%%%%%%%%%%%%%%%%%%%%%%%%%%%%%%%%%%%%%%%%%%%%%%%%%%%%

The true test of any theoretical framework is its ability to explain puzzles that resist conventional analysis. This section demonstrates how the Awareness framework diagnoses behavioral anomalies that standard economics cannot explain---and derives intervention strategies directly from the diagnosis.

\textbf{Key Result:} All nine cases yield predictions consistent with observed behavior when $U_{eff}$ is computed using calibrated awareness parameters.

%------------------------------------------------------------------------------
\subsubsection{The Diagnostic Method}
%------------------------------------------------------------------------------

Every behavioral puzzle can be analyzed through three steps:

\begin{enumerate}
    \item \textbf{Decompose}: What utility components are at stake? ($U_{INU}$, $U_{KNU}$, $U_{IDN}$)
    \item \textbf{Measure Awareness}: What is the $A$ value for each component at the Moment of Truth?
    \item \textbf{Compute $U_{eff}$}: Apply the transformation $U_{eff} = \sum_i A_i \cdot U_i$
\end{enumerate}

The sign of $U_{eff}$ predicts behavior; its magnitude indicates decision confidence.

%==============================================================================
% CASE-01: LABOR MARKET TRANSITIONS
%==============================================================================

\subsubsection{CASE-01: Labor Market Transitions}

\paragraph{The Puzzle.}
Coal miners reject retraining programs despite \euro{20,000}/year income gain. Observed rejection rate exceeds 70\%. Standard economics predicts acceptance.

\paragraph{Awareness Profile at Moment of Truth.}

\begin{equation}
\begin{array}{lclll}
\text{Component} & A & U & A \cdot U & \text{Source} \\
\midrule
A^{self}_{INU,F}(\text{salary gain}) & 0.3 & +1.0 & +0.30 & \text{SYS 2, delayed} \\
A^{self}_{KNU}(\text{community loss}) & 0.7 & -0.8 & -0.56 & \text{SYS 1, salient} \\
A^{self}_{IDN}(\text{identity threat}) & \mathbf{1.0} & -1.0 & -1.00 & \text{SYS 0/1, instant}
\end{array}
\end{equation}

\paragraph{Effective Utility Calculation.}
\begin{equation}
\boxed{U_{eff} = 0.30 - 0.56 - 1.00 = \mathbf{-1.26}}
\end{equation}

\paragraph{Prediction.}
$U_{eff} < 0 \Rightarrow$ \textbf{REJECT retraining}

\paragraph{Key Asymmetry.}
\begin{equation}
A^{self}_{IDN}(t^*) = 1.0 \quad \text{vs.} \quad A^{self}_{INU}(t^*) = 0.3
\end{equation}

Identity threat activates instantly (AWX-A36), while financial benefit requires deliberation. The 3:1 ratio means identity dominates---the \euro{20,000} is ``not seen'' at $t^*$.

\paragraph{Active Axioms.} AWX-A36 (Instant IDN), AWX-A34 (Hierarchy), AWX-A85 (Identity Stakes), AWX-A37 (No Instant KNU)

\paragraph{Active Corollaries.} COR-20 (Identity-Protective Cognition), COR-57 (Status Quo Bias), COR-43 (Loss Aversion)

\paragraph{Derived Interventions.}
\begin{itemize}
    \item \textbf{Target A36/A85}: Create identity bridge---``Energy Transition Leader'' not ``Ex-Miner''
    \item \textbf{Target A37}: Collective transition preserves community (KNU becomes visible)
    \item \textbf{Target A85}: Reduce identity stakes through gradual role evolution
\end{itemize}

%==============================================================================
% CASE-02: HEALTH AND PREVENTION
%==============================================================================

\subsubsection{CASE-02: Health and Prevention (Smoking)}

\paragraph{The Puzzle.}
Smokers continue despite knowing health consequences. Information campaigns show $<$10\% quit rates. Why does knowledge not translate to behavior?

\paragraph{Awareness Profile at Moment of Truth (Craving Moment).}

\begin{equation}
\begin{array}{lclll}
\text{Component} & A_0 & \text{Decay} & A_{eff} & U \\
\midrule
A^{self}_{INU,E,t=0}(\text{craving}) & 1.0 & - & \mathbf{1.0} & +1.0 \\
A^{self}_{INU,P,t=20y}(\text{cancer}) & 0.1 & e^{-0.1 \times 20} & 0.0135 & -5.0 \\
A^{self}_{IDN}(\text{smoker identity}) & 0.8 & - & 0.8 & +0.3
\end{array}
\end{equation}

\paragraph{Temporal Decay (AWX-A58).}
\begin{equation}
A_{future} = A_0 \cdot e^{-\lambda t} = 0.1 \cdot e^{-2.0} = 0.0135
\end{equation}

\paragraph{Effective Utility Calculation.}
\begin{equation}
U_{eff} = 1.0 \cdot (+1.0) + 0.0135 \cdot (-5.0) + 0.8 \cdot (+0.3)
\end{equation}
\begin{equation}
\boxed{U_{eff} = 1.0 - 0.068 + 0.24 = \mathbf{+1.172}}
\end{equation}

\paragraph{Prediction.}
$U_{eff} > 0 \Rightarrow$ \textbf{CONTINUE smoking}

\paragraph{Key Asymmetry.}
\begin{equation}
A_{now} = 1.0 \gg A_{future} = 0.0135 \quad \text{(74:1 ratio)}
\end{equation}

\paragraph{Why Information Fails.}
Information targets SYS 2, but SYS 0 dominates at Moment of Truth (AWX-A4: Hyper-Dissonance).

\paragraph{Active Axioms.} AWX-A4 (Instant Exception), AWX-A35 (Instant INU), AWX-A58 (Decay), AWX-A60 (Habituation), AWX-A84 (Asymmetric Processing)

\paragraph{Derived Interventions.}
\begin{itemize}
    \item \textbf{Target SYS 0}: Pharmacological interventions (reduce craving intensity)
    \item \textbf{Target A58}: Compress temporal distance (vivid future imagery)
    \item \textbf{Target A36}: Identity shift---``non-smoker'' before ``ex-smoker''
\end{itemize}

%==============================================================================
% CASE-03: FINANCIAL DECISIONS
%==============================================================================

\subsubsection{CASE-03: Financial Decisions (Undersaving)}

\paragraph{The Puzzle.}
Systematic undersaving despite knowing benefits. Median retirement savings $<$50\% of recommended level.

\paragraph{Awareness Profile.}

\begin{equation}
\begin{array}{lclll}
\text{Component} & A_0 & \text{Decay} & A_{eff} & U \\
\midrule
A^{self}_{INU,F,t=0}(\text{consume now}) & 0.9 & - & 0.9 & +1.0 \\
A^{self}_{INU,F,t=30y}(\text{retirement}) & 0.05 & e^{-0.12 \times 30} & 0.00137 & +2.0
\end{array}
\end{equation}

\paragraph{Extreme Temporal Decay.}
\begin{equation}
A_{retirement} = 0.05 \cdot e^{-3.6} = 0.05 \cdot 0.0273 = 0.00137
\end{equation}

\paragraph{Effective Utility Calculation.}
\begin{equation}
U_{eff} = 0.9 \cdot (+1.0) - 0.00137 \cdot (+2.0)
\end{equation}
\begin{equation}
\boxed{U_{eff} = 0.9 - 0.003 = \mathbf{+0.897}}
\end{equation}

\paragraph{Prediction.}
$U_{eff} > 0 \Rightarrow$ \textbf{UNDERSAVE} (consume now)

\paragraph{Key Asymmetry.}
\begin{equation}
A(t=0) = 0.9 \gg A(t=30y) = 0.00137 \quad \text{(657:1 ratio)}
\end{equation}

\paragraph{Why It Happens.}
AWX-A1b: Future self treated as ``other'' (community member, not self). AWX-A58: Exponential decay over 30 years effectively eliminates future awareness.

\paragraph{Derived Interventions.}
\begin{itemize}
    \item \textbf{Bypass $A$}: Auto-enrollment (default savings removes awareness requirement)
    \item \textbf{Target A1b}: Age progression imagery (connect present to future self)
    \item \textbf{Target A58}: Commitment devices (pre-commit before decay operates)
\end{itemize}

%==============================================================================
% CASE-04: CLIMATE ACTION
%==============================================================================

\subsubsection{CASE-04: Climate Action}

\paragraph{The Puzzle.}
Climate concern (85\%) vastly exceeds behavioral change (15\%). Largest awareness-action gap in behavioral economics.

\paragraph{Awareness Profile.}

\begin{equation}
\begin{array}{lclll}
\text{Component} & A & U & A \cdot U & \text{Note} \\
\midrule
A^{self}_{INU}(\text{cost now}) & 0.8 & -0.5 & -0.40 & \text{SYS 1} \\
A^{comm}_{KNU}(\text{collective}) & 0.1 & +1.0 & +0.10 & \text{SYS 2} \\
A^{society}_{KNU}(\text{global}) & 0.02 & +2.0 & +0.04 & \text{Abstract}
\end{array}
\end{equation}

\paragraph{AWX-A37: KNU is NEVER instant.}
\begin{equation}
A^{comm}_{KNU,instant} = 0 \quad \text{(by axiom definition)}
\end{equation}

\paragraph{Effective Utility Calculation.}
\begin{equation}
U_{eff} = -0.40 + 0.10 + 0.04
\end{equation}
\begin{equation}
\boxed{U_{eff} = \mathbf{-0.26}}
\end{equation}

\paragraph{Prediction.}
$U_{eff} < 0 \Rightarrow$ \textbf{NO climate action}

\paragraph{Key Asymmetry.}
\begin{equation}
A^{self}_{INU} = 0.8 \gg A^{comm}_{KNU} = 0.1 \quad \text{(8:1 ratio)}
\end{equation}

Individual costs are instant and salient; collective benefits are never instant (AWX-A37). This structural asymmetry explains the 70 percentage point gap.

\paragraph{Derived Interventions.}
\begin{itemize}
    \item \textbf{Target A37}: Make KNU visible and local (``Your neighborhood's solar panels'')
    \item \textbf{Target A20}: Social proof (``X\% of your neighbors switched'')
    \item \textbf{Target A55}: Institutional solutions (carbon pricing makes INU=0)
    \item \textbf{Target A40}: Align identity with KNU (``Climate leader'' identity)
\end{itemize}

%==============================================================================
% CASE-05: EDUCATION INVESTMENT
%==============================================================================

\subsubsection{CASE-05: Education Investment}

\paragraph{The Puzzle.}
First-generation students underinvest in education despite $>$15\% lifetime ROI.

\paragraph{Awareness Profile (Low SES).}

\begin{equation}
\begin{array}{lclll}
\text{Component} & A_0 & \text{SES Modifier} & A_{eff} & U \\
\midrule
A^{self}_{INU,t=0}(\text{opp. cost}) & 0.9 & - & 0.9 & -1.0 \\
A^{self}_{INU,t=10y}(\text{earnings}) & 0.2 & \times 0.3 & 0.06 & +3.0 \\
A^{self}_{IDN}(\text{identity barrier}) & 0.8 & - & 0.8 & -0.5
\end{array}
\end{equation}

\paragraph{SES Modulation (AWX-A75).}
Scarcity reduces SYS 2 capacity:
\begin{equation}
A_{SYS2,capacity} = 0.6 \times SES_{factor} = 0.6 \times 0.5 = 0.3
\end{equation}
\begin{equation}
A_{future,effective} = 0.2 \times 0.3 = 0.06
\end{equation}

\paragraph{Effective Utility Calculation.}
\begin{equation}
U_{eff} = 0.9 \cdot (-1.0) + 0.06 \cdot (+3.0) + 0.8 \cdot (-0.5)
\end{equation}
\begin{equation}
\boxed{U_{eff} = -0.9 + 0.18 - 0.4 = \mathbf{-1.12}}
\end{equation}

\paragraph{Prediction.}
$U_{eff} < 0 \Rightarrow$ \textbf{UNDERINVEST} in education

\paragraph{Key Mechanism.}
SES constraint reduces $A_{SYS2}$ from 0.2 to 0.06---a 70\% reduction in future awareness capacity. The 15\% ROI is ``invisible'' under cognitive load.

\paragraph{Derived Interventions.}
\begin{itemize}
    \item \textbf{Target A75}: Reduce immediate financial stress (grants, not loans)
    \item \textbf{Target A20}: Role models (``People like you who succeeded'')
    \item \textbf{Target A85}: Identity expansion (``Scholar'' compatible with background)
\end{itemize}

%==============================================================================
% CASE-06: POLITICAL POLARIZATION
%==============================================================================

\subsubsection{CASE-06: Political Polarization}

\paragraph{The Puzzle.}
Better-informed citizens show STRONGER bias. More information leads to more polarization (paradox).

\paragraph{Awareness Profile.}

\begin{equation}
\begin{array}{lcll}
A^{self}_{IDN}(\text{political identity}) & = & \mathbf{1.0} & \text{SYS 0/1, instant} \\
A^{self}_{INU,cog}(\text{policy analysis}) & = & 0.2 & \text{SYS 2}
\end{array}
\end{equation}

\paragraph{Belief Filter (AWX-A84).}
\begin{equation}
B_X = \begin{cases} 
1 & \text{if signal is belief-consistent} \\
1 - \lambda \cdot \left| \frac{\partial U_{IDN}}{\partial \hat{\theta}} \right| & \text{if belief-inconsistent}
\end{cases}
\end{equation}

With $\lambda_{politics} = 0.9$ (high identity stakes, AWX-A88):
\begin{equation}
B_{inconsistent} = 1 - 0.9 \times 1.0 = \mathbf{0.1}
\end{equation}

\paragraph{Result.}
90\% of inconsistent information is filtered $\Rightarrow$ \textbf{POLARIZATION}

\paragraph{Key Mechanism (AWX-A87).}
Higher cognitive capacity enables better \emph{motivated reasoning}, not better accuracy:
\begin{equation}
\text{Motivated Reasoning Power} = \text{Cognitive Capacity} \times \lambda_{identity} = 0.8 \times 0.9 = 0.72
\end{equation}

\paragraph{Derived Interventions.}
\begin{itemize}
    \item \textbf{Target A85}: Reduce identity stakes (depoliticize where possible)
    \item \textbf{Target A88}: Cross-cutting identities (multiple group memberships)
    \item \textbf{Target A89}: Break feedback loop through external accountability
\end{itemize}

%==============================================================================
% CASE-07: TECHNOLOGY ADOPTION
%==============================================================================

\subsubsection{CASE-07: Technology Adoption}

\paragraph{The Puzzle.}
Profitable technologies rejected by firms. Average adoption lag is 5-10 years despite clear ROI.

\paragraph{Awareness Profile.}

\begin{equation}
\begin{array}{lclll}
\text{Component} & A_0 & \text{K×U Penalty} & A_{eff} & U \\
\midrule
A_{current}(\text{status quo}) & 0.9 & - & 0.9 & +0.6 \\
A_{new}(\text{efficiency}) & 0.3 & \times (1-0.42) & 0.174 & +1.2 \\
A_{uncertainty}(\text{risk}) & 0.8 & - & 0.8 & -0.5
\end{array}
\end{equation}

\paragraph{Complexity × Uncertainty Interaction (AWX-A54).}
\begin{equation}
K \times U = 0.6 \times 0.7 = 0.42
\end{equation}
\begin{equation}
A_{new,adjusted} = 0.3 \times (1 - 0.42) = 0.174
\end{equation}

\paragraph{Utility Comparison.}
\begin{align}
U_{status\_quo} &= 0.9 \times 0.6 = \mathbf{0.54} \\
U_{change} &= 0.174 \times 1.2 + 0.8 \times (-0.5) = 0.209 - 0.4 = \mathbf{-0.191}
\end{align}

\paragraph{Prediction.}
$U_{change} < U_{status\_quo} \Rightarrow$ \textbf{REJECT/DELAY adoption}

\paragraph{Key Asymmetry.}
Status quo requires $A = 0$ (habituated); change requires $A > 0$ plus overcoming $K \times U$ penalty.

\paragraph{Derived Interventions.}
\begin{itemize}
    \item \textbf{Target A46}: Reduce complexity (phased rollout, training)
    \item \textbf{Target A50}: Reduce uncertainty (pilot programs, guarantees)
    \item \textbf{Target A54}: Address $K$ and $U$ separately, never together
\end{itemize}

%==============================================================================
% CASE-08: ORGANIZATIONAL CHANGE
%==============================================================================

\subsubsection{CASE-08: Organizational Change}

\paragraph{The Puzzle.}
70\% of organizational change initiatives fail despite clear strategic rationale.

\paragraph{Awareness Profile.}

\begin{equation}
\begin{array}{lcll}
A^{mgmt}_{KNU}(\text{strategic benefit}) & = & 0.8 & \text{SYS 2} \\
A^{employee}_{INU}(\text{disruption}) & = & 0.9 & \text{SYS 1} \\
A^{employee}_{IDN}(\text{role identity}) & = & 0.7 & \text{SYS 1}
\end{array}
\end{equation}

\paragraph{Collective Dynamics (AWX-A70, A78, A80).}

Base resistance rate: 60\%

Cascade multiplier (AWX-A78): 1.3

\begin{equation}
Resistance_{cascaded} = 0.6 \times 1.3 = \mathbf{0.78}
\end{equation}

\paragraph{Blockade Condition (AWX-A80).}
\begin{equation}
Resistance_{cascaded} = 0.78 > \theta_{blockade} = 0.5
\end{equation}

\paragraph{Prediction.}
Blockade triggered $\Rightarrow$ \textbf{CHANGE FAILS}

\paragraph{Formula.}
\begin{equation}
A_{Group,resistance} = \sum_i w_i A_i + \text{Cascade} > \theta_{blockade} \Rightarrow \text{Change fails}
\end{equation}

\paragraph{Derived Interventions.}
\begin{itemize}
    \item \textbf{Target A70}: Build coalition of early adopters (prevent collective override)
    \item \textbf{Target A36}: Address identity first---``Your expertise transfers''
    \item \textbf{Target A78}: Staged rollout to prevent cascade
\end{itemize}

%==============================================================================
% CASE-09: CONSUMER BEHAVIOR
%==============================================================================

\subsubsection{CASE-09: Consumer Behavior (Sustainable Consumption)}

\paragraph{The Puzzle.}
65\% intend sustainable purchase, 25\% execute. Attitude-behavior gap = 40 percentage points.

\paragraph{Awareness: Survey vs. Store.}

\begin{equation}
\begin{array}{lcll}
\text{Context} & \text{Component} & A & U \\
\midrule
\text{Survey} & A_{KNU}(\text{values}) & 0.7 & +0.8 \\
\text{Store} & A_{INU}(\text{price}) & 0.95 & -0.3 \\
\text{Store} & A_{KNU}(\text{environment}) & 0.15 & +0.8
\end{array}
\end{equation}

\paragraph{AWX-A2/A5: Latent $\neq$ Manifest.}
\begin{equation}
A_{survey} = 0.7 \neq A_{store}(t^*) = 0.15
\end{equation}

\paragraph{Effective Utility at Point of Sale.}
\begin{equation}
U_{store} = 0.95 \times (-0.3) + 0.15 \times (+0.8)
\end{equation}
\begin{equation}
\boxed{U_{store} = -0.285 + 0.12 = \mathbf{-0.165}}
\end{equation}

\paragraph{Prediction.}
$U_{store} < 0 \Rightarrow$ \textbf{BUY CONVENTIONAL}

\paragraph{Key Asymmetry.}
\begin{equation}
A_{survey} = 0.7 \quad \text{vs.} \quad A_{store} = 0.15 \quad \text{(4.7:1 ratio)}
\end{equation}

Surveys measure general awareness; behavior requires situational awareness at $t^*$.

\paragraph{Derived Interventions.}
\begin{itemize}
    \item \textbf{Target A9}: Point-of-sale cues (make sustainability salient at $t^*$)
    \item \textbf{Bypass $A$}: Default to sustainable option
    \item \textbf{Target A20}: Social visibility (``X\% chose sustainable'')
\end{itemize}

%==============================================================================
% SUMMARY TABLE
%==============================================================================

\subsubsection{Summary: All Nine Cases}

\begin{table}[htbp]
\centering
\caption{Application Cases: Computed Outcomes}
\label{tab:case-computed-summary}
\small
\begin{tabular}{llccc}
\toprule
\textbf{Case} & \textbf{Domain} & $\mathbf{U_{eff}}$ & \textbf{Prediction} & \textbf{Observed} \\
\midrule
CASE-01 & Labor Transitions & $-1.26$ & Reject & $>$70\% reject \\
CASE-02 & Health/Smoking & $+1.17$ & Continue & $<$10\% quit \\
CASE-03 & Financial/Saving & $+0.90$ & Undersave & $<$50\% adequate \\
CASE-04 & Climate Action & $-0.26$ & No action & 15\% act \\
CASE-05 & Education & $-1.12$ & Underinvest & Gap persists \\
CASE-06 & Polarization & $B=0.1$ & Filter info & Paradox confirmed \\
CASE-07 & Tech Adoption & $-0.19$ & Delay & 5-10yr lag \\
CASE-08 & Org. Change & 78\%$>$50\% & Fails & 70\% fail \\
CASE-09 & Sust. Consumption & $-0.17$ & Conventional & 40pp gap \\
\bottomrule
\end{tabular}
\end{table}

\textbf{Result: 9/9 cases correctly predicted by the AWX framework.}

%==============================================================================
% CROSS-DOMAIN PATTERNS
%==============================================================================

\subsubsection{Cross-Domain Patterns}

Five mechanisms recur across all nine cases:

\begin{enumerate}
    \item \textbf{Temporal Decay (AWX-A58)}: 
    \begin{equation}
    A(t) = A_0 \cdot e^{-\lambda t} \Rightarrow \text{Future awareness } \to 0
    \end{equation}
    CASE-02: 74:1 now/future ratio. CASE-03: 657:1 ratio over 30 years.
    
    \item \textbf{Instant-Asymmetry (AWX-A36/A37)}: 
    \begin{equation}
    A_{IDN}^{instant} = 1, \quad A_{KNU}^{instant} = 0 \text{ (always)}
    \end{equation}
    Identity dominates in CASE-01, CASE-06. KNU disadvantaged in CASE-04, CASE-09.
    
    \item \textbf{Belief Filter (AWX-A84)}:
    \begin{equation}
    B_{inconsistent} = 1 - \lambda \cdot |U'_{IDN}| \approx 0.1 \text{ for political beliefs}
    \end{equation}
    Explains CASE-06 polarization paradox.
    
    \item \textbf{K×U Interaction (AWX-A54)}:
    \begin{equation}
    A_{new} \cdot (1 - K \times U) \Rightarrow 42\% \text{ penalty in CASE-07}
    \end{equation}
    Complexity combined with uncertainty is fatal for adoption.
    
    \item \textbf{Collective Cascade (AWX-A78/A80)}:
    \begin{equation}
    Resistance_{cascaded} > \theta_{blockade} \Rightarrow \text{Organizational change fails}
    \end{equation}
    Individual resistance amplifies to collective blockade in CASE-08.
\end{enumerate}

\subsubsection{The Power of Formal Diagnosis}

Each case demonstrates the same logic:
\begin{enumerate}
    \item \textbf{Puzzle}: Standard economics cannot explain the behavior
    \item \textbf{Diagnosis}: Awareness asymmetry identified via specific axioms
    \item \textbf{Computation}: $U_{eff}$ calculated with calibrated parameters
    \item \textbf{Verification}: Sign of $U_{eff}$ matches observed behavior
    \item \textbf{Intervention}: Target the specific axioms creating the asymmetry
\end{enumerate}

The interventions are not ad hoc but \emph{derived} from the formal structure. This is the practical power of axiomatic behavioral economics.

\begin{cross-reference}
\textbf{Formal Specification:} Complete awareness profiles, all active axioms and corollaries, and formal expressions for each case are documented in Appendix~\ref{app:bcm-axiom-formalization}, Section ``Application Cases'' (CASE-01 through CASE-09).
\end{cross-reference}

%%%%%%%%%%%%%%%%%%%%%%%%%%%%%%%%%%%%%%%%%%%%%%%%%%%%%%%%%%%%%%%%%%%%%%%%%%%%%%%
% END OF SECTION 11.16
%%%%%%%%%%%%%%%%%%%%%%%%%%%%%%%%%%%%%%%%%%%%%%%%%%%%%%%%%%%%%%%%%%%%%%%%%%%%%%%
